% ch3 D.3 铁磁自旋波 (原书页 164–174 / PDF p179–p189)

% 衔接说明：%--D2-TAIL-- 标记之下是前一小节 D.2 末句的续文（原书 p.164 顶部、
% 小节标题 “3. Ferromagnetic Spin Waves” 之前的一段），装配时应移至承接原书
% p.163 的 D.2 文件（ch3_D2.tex，其末句中断于 “The lower iron-series oxides,”）
% 末尾拼接。Figure 55 印于 p.164 页首、该段续文之前，故 FIG 标记一并置于
% 同一区域内。
%--D2-TAIL--
%JOIN%
以及后面几个过渡系列——V、Ru 等——的氧化物则是金属型的，尽管其中电子的迁移率相当低，只有通常金属中的 $10^{-3}$ 左右。有些氧化物表现出一种热致转变，有人提出，它也许在本质上是一种由金属型到磁性的转变 (9)。
%FIG{55}: 图 55　（原书图注仅作 “Figure 55”）p.163 自洽方程 $(\Delta E)^{3} - U(\Delta E)^{2} = -2UZ|b_{jk}|^{2}$ 的图解：纵轴为 $\Delta E/U$，横轴为 $2Z|b/U|^{2}$；曲线自原点出发，在虚线标出的 $\Delta E/U = 2/3$ 处达到极大，并于横坐标 $4/27$ 处陡降为零

\subsection{铁磁自旋波}

现在让我们继续讨论元激发（elementary excitations），先讨论铁磁情形 (96)。有效哈密顿量是

\begin{equation*}
\mathcal{H} = -\frac{1}{4} \sum_{jk} J_{jk}\, \boldsymbol{\sigma}_{j}\cdot\boldsymbol{\sigma}_{k}
\end{equation*}

（这里假定 $J_{jk} > 0$。）

自旋矢量与谐振子共有的一个最重要的性质是：它们的经典运动方程与量子运动方程完全相同。设 $\boldsymbol{\sigma}$ 与一个场 $\mathbf{H}$ 相乘后进入哈密顿量：$\mathcal{H} = -\frac{1}{2}\gamma\hbar\,\mathbf{H}\cdot\boldsymbol{\sigma}$。于是，当我们在 Heisenberg 表象中计算算符 $\boldsymbol{\sigma}$ 的运动方程时：

\begin{equation*}
i\hbar\,\frac{d\boldsymbol{\sigma}}{dt} = \left[\boldsymbol{\sigma},\, \mathcal{H}\right]
\end{equation*}

便得到

\begin{equation*}
i\hbar\,\frac{d\boldsymbol{\sigma}}{dt} = -\frac{\gamma\hbar}{2}\left[\boldsymbol{\sigma},\boldsymbol{\sigma}\right]\cdot\mathbf{H}
\end{equation*}

（这一记号的意思是：$[\boldsymbol{\sigma},\boldsymbol{\sigma}]$ 表示由分量 $[\sigma_{i},\sigma_{j}]$ 组成的并矢（dyadic）。）不同分量的对易子概括在关系式

\begin{equation*}
\boldsymbol{\sigma}\times\boldsymbol{\sigma} = 2i\boldsymbol{\sigma}
\end{equation*}

或者

\begin{equation*}
\sigma_{i}\sigma_{j} - \sigma_{j}\sigma_{i} = 2i\,\epsilon_{ijk}\,\sigma_{k}
\end{equation*}

之中，于是

\begin{equation*}
\left[\left[\boldsymbol{\sigma},\boldsymbol{\sigma}\right]\mathbf{H}\right]_{k} = 2i\sum_{j} H_{j}\,\epsilon_{kjl}\,\sigma_{l} = 2i\,\mathbf{H}\times\boldsymbol{\sigma}
\end{equation*}

因此

\begin{equation*}
\frac{d\boldsymbol{\sigma}}{dt} = \gamma\left(\boldsymbol{\sigma}\times\mathbf{H}\right)
\end{equation*}

这正是处于磁场 $\mathbf{H}$ 中、回磁比（gyromagnetic ratio）为 $\gamma$ 的自旋 $\mathbf{S} = \boldsymbol{\sigma}/2$ 的经典运动方程

\begin{equation*}
\frac{d\mathbf{S}}{dt} = \gamma\left(\mathbf{S}\times\mathbf{H}\right)
\end{equation*}

$\mathbf{H}$ 可以是一个 c 数（c-number），也可以是任何与所考虑的自旋 $\boldsymbol{\sigma}$ 对易的算符。

自旋 $\boldsymbol{\sigma}_{j}$ 所感受到的场 $\mathbf{H}$ 是 $\frac{1}{2\gamma\hbar}\sum_{k} J_{jk}\boldsymbol{\sigma}_{k}$。为方便起见，可以再附加一个小外场来引入一个优先方向 $\hat{\mathbf{z}}$，这在哈密顿量中导致一个 $\hbar\gamma H\sigma_{z}$ 项，于是我们便有运动方程

\begin{equation*}
\frac{d\boldsymbol{\sigma}_{j}}{dt} = \frac{1}{2\hbar}\sum_{k} J_{jk}\,\boldsymbol{\sigma}_{j}\times\boldsymbol{\sigma}_{k} + \gamma H\left(\boldsymbol{\sigma}_{j}\times\hat{\mathbf{z}}\right) \tag{24}
\end{equation*}

（译注：式 (24) 左端原书印作 $d\boldsymbol{\sigma}_{z}/dt$；按右端当为 $\boldsymbol{\sigma}_{j}$，此处径改。）

这些方程可以用来得到自旋波的运动方程，它们与用玻色子理论导出的方程完全等价；而且利用自旋和运动方程，许多性质可以直截了当地算出，根本不必经过赝玻色子近似（pseudo-boson approximation）。遗憾的是，热学的和能量的性质用这种方法却不那么容易处理，尽管近代的 Green 函数理论已在这个方向上取得了相当大的进展。尽管如此，记住下面这一点仍然十分重要：在自旋理论中，正如在谐振子理论中一样，大多数经典结果根本不是近似，而且人们的经典直觉可以具有极大的价值。

让我们把运动方程 (24) 明写出来：

\begin{align*}
\frac{d\sigma_{jx}}{dt} &= \frac{1}{2\hbar}\sum_{k} J_{jk}\left(\sigma_{jy}\sigma_{kz} - \sigma_{ky}\sigma_{jz}\right) + \gamma\sigma_{jy}H \\
\frac{d\sigma_{jy}}{dt} &= \frac{1}{2\hbar}\sum_{k} J_{jk}\left(\sigma_{jx}\sigma_{kz} - \sigma_{kx}\sigma_{jz}\right) - \gamma\sigma_{jx}H
\end{align*}

把第二式乘以 $-i$ 加到第一式上，我们得到

\begin{align*}
\frac{d\sigma^{-}_{j}}{dt} &= \frac{1}{2}\,\frac{d}{dt}\left(\sigma_{jx} - i\sigma_{jy}\right) \\
&= \frac{i}{2\hbar}\sum_{k} J_{jk}\left(\sigma^{-}_{j}\sigma_{kz} - \sigma^{-}_{k}\sigma_{jz}\right) + i\gamma\sigma^{-}_{j}H
\end{align*}

只有 $\sigma^{-}$ 出现这一事实表明，自旋波是圆偏振的。既然我们感兴趣的是把各种性质——热的与动力学的——都计算出来，那就让我们作玻色子变换

\begin{equation*}
\sigma^{-}_{j} = b^{\dagger}_{j}\left(1-n_{j}\right) \qquad\qquad \sigma_{jz} = 1 - 2n_{j}
\end{equation*}

我们把符号反了过来，原因是基态对应于 $\sigma_{z} = +1$，而一个自旋波所涉及的是把一个自旋转\emph{向下}。$b$ 的线性项是

\begin{equation*}
\frac{db^{\dagger}_{j}}{dt} = i\gamma H b^{\dagger}_{j} + \frac{i}{2\hbar}\sum_{k} J_{jk}\left(b^{\dagger}_{j} - b^{\dagger}_{k}\right)
\end{equation*}

当我们变换到自旋波变量 $\beta_{k}$ 时，我们得到

\begin{align*}
\frac{d\beta^{\dagger}_{k}}{dt} &= i\omega_{k}\beta^{\dagger}_{k} \\
\omega_{k} &= \gamma H + \frac{1}{2\hbar}\sum_{j} J_{0j}\left(1 - \cos\mathbf{k}\cdot\mathbf{j}\right) \tag{25}
\end{align*}

在没有磁场时，自旋波频率随 $k^{2}$ 从 $k=0$ 处的零值增大。除外场贡献之外频率在 $k=0$ 处必须为零，这一事实是下述恒等式的后果，该恒等式来自普遍的运动方程：

\begin{equation*}
\frac{d}{dt}\sum_{j}\sigma^{-}_{j} = \frac{i}{2\hbar}\sum_{jk} J_{jk}\left(\sigma^{-}_{j}\sigma_{kz} - \sigma^{-}_{k}\sigma_{jz}\right) \equiv 0
\end{equation*}

而且事实上，$S_{\mathrm{o}} = \sum_{j}\boldsymbol{\sigma}_{j}$，即总自旋矢量本身，是一个运动常数。$S^{-}$ 是 $k=0$ 自旋波算符的 $\sqrt{N}$ 倍。

这是我们第一次碰到集体激发的最重要普遍法则中的一条，它与哈密顿量中对称性的存在有关。交换哈密顿量 (23)，以及它所由之导出的原始哈密顿量，在自旋空间中的转动操作下不变。这是因为 Coulomb 相互作用和动能都不以任何方式涉及电子的自旋；自旋只是通过不相容原理、或者通过费米统计更完备的表述——对易关系——而进入问题。但对易关系本身在自旋转动（当然是同时转动所有自旋）之下也是不变的：

\begin{equation*}
\left[\Psi_{\sigma}(r),\ \Psi^{\dagger}_{\sigma}(r')\right] = \delta_{\sigma\sigma'}\,\delta(r-r')
\end{equation*}

在对应于自旋空间中一个转动的变换之后依然成立：

\begin{equation*}
\Psi_{+} \rightarrow \cos\frac{\theta}{2}\, e^{-i(\Phi/2)}\,\Psi^{'}_{+} + \sin\frac{\theta}{2}\, e^{i(\Phi/2)}\,\Psi^{'}_{-}
\end{equation*}

于是，当我们沿 $\theta,\Phi$ 方向而不是沿 $z$ 方向量子化时，$\Psi^{'}_{+}$ 就是波函数的 $+$ 分量（图 56）。
%FIG{56}: 图 56　（原书图注仅作 “Figure 56”）坐标系示意图：$z'$ 轴相对 $z$ 轴的极角为 $\theta$，绕 $z$ 轴的方位角为 $\phi$

哈密顿量的每一个不变性都对应一个运动常数。对应于时间平移不变性，该常数是 $E$；对应于空间平移不变性，是 $\mathbf{P}$；对应于角向转动不变性，则是总角动量 $\mathbf{L}$。它们就是所谓的对称变换的“生成元”（generators）；例如，$L_{z} = i(\partial/\partial\Phi)$，而当 $\partial\mathcal{H}/\partial\Phi = 0$ 时，

\begin{equation*}
\frac{\partial}{\partial\Phi}\,\mathcal{H} - \mathcal{H}\,\frac{\partial}{\partial\Phi} = 0 \qquad \text{于是} \qquad \left[\mathcal{H},\ L_{z}\right] = 0
\end{equation*}

于是，自旋空间中的转动不变性意味着 $[\mathcal{H},\mathbf{S}] = 0$，即 $\mathbf{S} = \text{const.}$。实际上这并不是一个严格的对称性，因为电磁力中磁性的部分确实依赖于 $\boldsymbol{\sigma}$；它只在 $v/c \rightarrow 0$ 的极限下成立，因此，比如说，自旋-轨道耦合就破坏它。结果，真实的铁磁体具有所谓的“晶体各向异性能”（crystal anisotropy energy），它依赖于 $\mathbf{S}$ 在空间中的方向，但其大小几乎总是比交换能小得多——交换能来自 Coulomb 能，并且在基本上决定了磁性性质。

这一近似对称性保证了 $k=0$ 自旋波的频率应当接近于零。由 $k=0$ 自旋波的激发所得出的激发态是什么？原来的铁磁态恰好含有 $N$ 个电子，每个都处于自旋朝上的态，因此总自旋为 $S = N/2$（$\sigma_{\mathrm{tot}} = N$），精确地指向 $z$ 方向。由于哈密顿量的转动不变性，$z$ 方向当然是完全任意的；我们当初未尝不可以沿 $x$ 或 $y$ 方向、或沿任何其他方向量子化，只要我们愿意。与此相应，基态存在简并；坚持取 $z$ 为量子化方向，我们本可以选

\begin{equation*}
S_{z} = -\frac{N}{2},\ -\frac{N}{2} + 1,\ \cdots,\ 0,\ 1,\ \cdots,\ \frac{N}{2} - 1,\ \frac{N}{2}
\end{equation*}

而这些严格本征态中的每一个都与原来的态具有完全相同的能量。对这些态取适当的线性组合，我们就可以构成一个 $S_{z'} = N/2$ 的态，其 $z'$ 指向任意方向 $\theta,\Phi$。应当取什么样的线性组合，可以由相应的表示系数 $D^{N/2}_{N/2,\,N/2}(\theta,\Phi)$ 算出；或者用这样一个技巧：把（这 $N$ 个电子的）态看作由乘积

\begin{equation*}
\alpha(1)\,\alpha(2)\cdots\alpha(N)
\end{equation*}

构成［$\alpha(i)$ 是与 $\sigma_{j}$ 相对应的波函数的自旋朝上分量］，并利用变换

\begin{equation*}
\alpha \rightarrow \alpha^{'}\cos\frac{\theta}{2}\, e^{i(\Phi/2)} + \beta^{'}\sin\frac{\theta}{2}\, e^{-i(\Phi/2)}
\end{equation*}

显然，从 $S_{z} = N/2$、$S = N/2$ 的态出发，逐次激发 $k=0$ 自旋波所得到的态，是 $S = N/2$ 但 $S_{z} = (N/2) - 1$、$(N/2) - 2$ 等等的态。这是因为

\begin{equation*}
\beta^{\dagger}_{k=0} = \frac{1}{\sqrt{N}}\sum_{j}\sigma^{-}_{j} = \frac{1}{\sqrt{N}}\,S^{-}
\end{equation*}

而 $S^{-}$ 就是“下降算符”（lowering operator），它使 $S_{z}$ 减小一个单位而不改变 $|S|$。这样我们看到，基态的转动不变性与 $k=0$ 模的零频率是等价的。

即使存在磁场，由 $k=0$ 模激发所到达的本征态仍是同样的这些态，因为在微扰 $\hbar\gamma H S_{z}$ 之下 $S_{z}$ 和 $|S|$ 仍然是好量子数。这时 $k=0$ 自旋波的频率不是零而恰好是 $\gamma H$。注意，关于均匀模频率的这些论断，在不存在自旋-轨道力的前提下，对哈密顿量的所有态都成立；这是显而易见的，因为 $S^{-}$ 的作用不过是把任何一个态转动一下。这就是那条著名的定理：“单是交换并不会展宽磁共振频率。”换一种说法：$k=0$ 自旋波与 $k\neq 0$ 的自旋波不以任何方式发生相互作用。

那么，$k\neq 0$ 的自旋波又如何？有一个显而易见的假设（而且事实证明确实如此）：当只考虑交换力时，自旋波的频率将连续地衔接到 $k=0$ 模的频率——也就是说，当 $k \rightarrow 0$ 时 $\omega \rightarrow 0$。

鉴于 $k \rightarrow 0$ 极限在其他问题中的种种复杂性，这一点并非显而易见，它是交换力短程性质的一个结果。例如，当我们把长程磁偶极力计入时，$k=0$ 模的频率就要依赖于样品的形状，因为表面磁极引起的退磁效应依赖于样品形状。而且纵向与横向自旋波的频率各不相同，所以当 $k \rightarrow 0$ 时根本谈不上连续性。与此相反，交换力基本上是短程的，因为除去磁偶极效应和自旋-轨道效应之外，自旋矩的重新分布在长距离上不会以任何可察觉的方式影响电荷分布。

这就是我想强调的重要联系，因为它在集体激发和相变的普遍理论中起着至关重要的作用。我们从一个哈密顿量

\begin{equation*}
\frac{1}{4} \sum_{jk} J_{jk}\,\boldsymbol{\sigma}_{j}\cdot\boldsymbol{\sigma}_{k}
\end{equation*}

（译注：此式原书未印负号；对照本节开头 $\mathcal{H} = -\frac{1}{4}\sum_{jk}J_{jk}\boldsymbol{\sigma}_{j}\cdot\boldsymbol{\sigma}_{k}$，疑为漏印，照录原书。）

出发，它在自旋空间中具有完全的转动对称性。尽管有这种转动对称性，我们所假定的那个最低本征态——它在某些特殊条件下、即所有 $J_{jk} < 0$ 时\emph{确}实是最低本征态——却并不具有完全的转动对称性。只有单态（$S_{\mathrm{tot}} = 0$）才具有。因此，最低本征态必定与大量其他态简并：这些态指向不同的方向，或者（如果我们愿意这样看）沿同一方向具有不同的 $S_{z}$。

于是我们得出结论：至少存在一种形式的元激发，即对应于 $k=0$ 的那一种，它给出整个系统在自旋空间中的一个转动，并且其频率必为 $\omega = 0$。在再加一条物理条件——不存在长程力——之下，这就意味着元激发能谱中有整整一支，当 $k \rightarrow 0$ 时 $\omega \rightarrow 0$。

如今，自旋波已经成为一种和声子同样实实在在的实验现象。人们已用中子非弹性散射、用薄膜中磁共振的直接观测、用与声子谱的耦合以及若干其他较为间接的方法对它进行过研究。不过，它最为人熟知的性质，是它对磁化强度和比热的温度依赖关系的影响。频率为 $\hbar\omega = Ak^{2}$ 的一支玻色型激发具有数密度

\begin{equation*}
N = 4\pi\int_{0}^{k_{\max}} \frac{k^{2}\,dk}{e^{\beta Ak^{2}} - 1} \qquad\qquad \beta = \frac{1}{k_{\mathrm{B}}T}
\end{equation*}

而当 $Ak_{\max}^{2} \gg k_{\mathrm{B}}T$ 时——即在低温下——这显然变为

\begin{align*}
N \approx 4\pi\int_{0}^{\infty} \frac{k^{2}\,dk}{e^{\beta Ak^{2}} - 1} &= \frac{4\pi}{(\beta A)^{3/2}}\int_{0}^{\infty} \frac{x^{2}\,dx}{e^{x} - 1} \\
&= \frac{4\pi\,(k_{\mathrm{B}}T)^{3/2}}{A^{3/2}} \times \text{const.}
\end{align*}

现在我们注意到，每个自旋波所对应的反转自旋数恰好是 1，因此 $S_{z}$ 的改变就正比于 $T^{3/2}$——这正是 Bloch 理论中著名的温度变化关系。Dyson (81) 已对这一变化关系计算了修正项：来自对 $\hbar\omega = Ak^{2}$ 的偏离的有 $T^{5/2}$ 和 $T^{7/2}$ 项，而第一个自旋波相互作用项则是 $T^{4}$ 量级的项，完全可以忽略。我们起初也许会以为 $A(T) \propto \text{const.} - N$，从而预期有一个 $T^{3}$ 项；但是长波长自旋波之间的相互作用非常小，其原因基本上是与 $k=0$ 波的连续衔接——后者与其他自旋波没有任何相互作用。人们发现

\begin{equation*}
\omega \propto k^{2}\left[\text{const.} - \int \left[k'^{2}N(k')\right] k'^{2}dk'\right] \propto k^{2}\left[\text{const.} - T^{5/2}\right]
\end{equation*}

其中 $N(k)$ 是波矢为 $\mathbf{k}$ 的自旋波的平均数目。相应地，修正项就是 $T^{3/2}T^{5/2} = T^{4}$。自旋波比热同样 $\propto T^{3/2}$：能量为

\begin{equation*}
U = 4\pi A\int \frac{k^{2}}{e^{\beta Ak^{2}} - 1}\, k^{2}dk \propto T^{5/2} \qquad\qquad C = \frac{\partial U}{\partial T} \propto T^{3/2}
\end{equation*}

这些积分在 $k=0$ 附近的行为说明了一个重要的论点。在任一有限温度下、靠近 $k=0$ 处，自旋波密度 $N(k)$ 就是经典值

\begin{equation*}
N(k) = \frac{k_{\mathrm{B}}T}{\hbar\omega(k)} = \frac{k_{\mathrm{B}}T}{Ak^{2}}
\end{equation*}

请注意，假如系统只有两维，从而态密度 $\propto k\,dk$，那么磁化偏离的积分将在 $k \rightarrow 0$ 处对数发散。不过，任何小的外场或自旋-轨道耦合都可以使 $k=0$ 自旋波的频率变为有限，从而使积分收敛，其形式大致为

\begin{equation*}
T\ln\left(\frac{k_{\mathrm{B}}T}{\text{截断能量}}\right)
\end{equation*}

这种二维发散也是具有短程力的有序系统的一个非常普遍的性质。例如，人们相信，如果没有某种有序的支撑，二维晶体将不是热力学稳定的。然而，对于二维情形不存在相变这一点，看来实际上并没有证明 (98)。

现在我想把这些结果与两种不同类型的经典理论的结果联系起来。第一种是纯经典自旋矢量所对应的同一个问题。假定代替 Pauli 自旋算符 $\boldsymbol{\sigma}_{j}$，我们有的是经典自旋矢量 $\mathbf{S}_{j}$：

\begin{equation*}
\mathcal{H} = -\sum_{jk} J_{jk}\,\mathbf{S}_{j}\cdot\mathbf{S}_{k}
\end{equation*}

我们定义 $\mathbf{S}_{j} = \mathbf{J}_{j}/\hbar$，其中 $\mathbf{J}_{j}$ 是真实的经典角动量。这样，经典交换参量便是 $J_{jk}/\hbar^{2}$。我们指出这一点，是因为我们人为地使用 $\mathbf{S}$ 而不是 $\mathbf{J}$，可能会在实际上是经典的方程中引进 $\hbar$ 因子。

这个哈密顿量的最低态，应当是所有经典自旋矢量都完全沿同一方向排列的态。我们可以把这一能量在平衡位置附近展开，从而用这种阵列的经典进动模导出一个近似哈密顿量：

\begin{equation*}
\mathbf{S}_{j}\cdot\mathbf{S}_{k} = S_{zj}S_{zk} + S_{xj}S_{xk} + S_{yj}S_{yk}
\end{equation*}

\begin{equation*}
S_{zj}^{2} = S^{2} - S_{xj}^{2} - S_{yj}^{2}
\end{equation*}

于是

\begin{equation*}
S_{zj} \simeq S - \frac{1}{2S}\left(S_{xj}^{2} + S_{yj}^{2}\right)
\end{equation*}

因为在平衡位置附近第二项很小。这样便有

\begin{align*}
-\mathcal{H} &= -NZJS^{2} + \sum_{jk} J\left[\left(S_{xj} - S_{xk}\right)^{2} + \left(S_{yj} - S_{yk}\right)^{2}\right] \\
&= -NZJS^{2} + \sum_{\substack{jk \\ \text{近邻}}} \frac{J}{2}\left[\left(S^{+}_{j} - S^{+}_{k}\right)\left(S^{-}_{j} - S^{-}_{k}\right) + \left(S^{-}_{j} - S^{-}_{k}\right)\left(S^{+}_{j} - S^{+}_{k}\right)\right]
\end{align*}

（译注：首行左端原书印作 $-\mathcal{H}$；按上页 $\mathcal{H} = -\sum_{jk}J_{jk}\mathbf{S}_{j}\cdot\mathbf{S}_{k}$ 推算当为 $\mathcal{H}$，与下文式 (26) 对照亦可见负号为误印，照录原书。）

这里我们假定了 $J_{jk}$ 除原子与其 $Z$ 个近邻之间外处处为零，这是一种非本质的简化。$S^{-}$ 与 $S^{+}$ 的定义与量子力学自旋的情形相同。我们还可以再前进一步，注意到无论经典地还是量子力学地，自旋的小振荡 $\hbar^{1/2}(\delta S_{x}/\sqrt{S})$ 与 $\hbar^{1/2}(\delta S_{y}/\sqrt{S})$ 都服从正则共轭变量的 Poisson 括号关系。更简单地说，我们指出，自旋的经典运动方程与量子运动方程是等同的，因此系统的经典模同样是进动波：

\begin{equation*}
S^{-}_{j}(t) = e^{i\omega_{\lambda} t}\, e^{i\boldsymbol{\lambda}\cdot\mathbf{j}}
\end{equation*}

也就是说，我们可以导出一组简正模（normal modes）

\begin{align*}
S^{-}_{\lambda} &= \frac{1}{\sqrt{NS}}\sum_{j} e^{i\boldsymbol{\lambda}\cdot\mathbf{j}}\, S^{-}_{j} \\
S^{-}_{\lambda} &\propto e^{i\omega_{\lambda} t}
\end{align*}

$\omega_{\lambda}$ 由下式给出：

\begin{equation*}
\omega_{\lambda} = \frac{S}{\hbar}\sum_{j} J_{0j}\left(1 - \cos\boldsymbol{\lambda}\cdot\mathbf{j}\right)
\end{equation*}

若在 (25) 中代入 $S = 1/2$，此式便与 (25) 相一致；而如果我们代入 $J = \hbar S$，并意识到 $J_{0j}$ 是经典交换参量的 $\hbar^{2}$ 倍，它也还是一个纯经典的公式。哈密顿量约化为

\begin{equation*}
\mathcal{H} = -NZJS^{2} + \sum_{\lambda} \frac{\omega_{\lambda}}{2}\left(S^{+}_{\lambda}S^{-}_{\lambda} + S^{-}_{\lambda}S^{+}_{\lambda}\right) \tag{26}
\end{equation*}

当然，经典地看次序无关紧要；但当我们把系统量子化时，就必须像通常那样选取这个对称化了的厄米组合。

关于零点运动和零点能量的情形起初是令人困惑的。量化的经典 Hamilton 量与相应的量子结果之间好像有出入，因为量子力学中基态能量恰好是 $-NZJS^{2}$，这里没有为对易子项带来的修正留有余地——这一修正必定来自把 $S^{+}S^{-} + S^{-}S^{+}$ 项重新编序并加以修正：

\begin{align*}
\frac{1}{2}&\left\{\left(S^{x}_{\lambda} + iS^{j}_{\lambda}\right)\left(S^{x}_{\lambda} - iS^{y}_{\lambda}\right) - \left(S^{x}_{\lambda} - iS^{y}_{\lambda}\right)\left(S^{x}_{\lambda} + iS^{y}_{\lambda}\right)\right\} \\
&= -i\left(S^{x}_{\lambda}S^{y}_{\lambda} - S^{y}_{\lambda}S^{x}_{\lambda}\right) = i\frac{1}{NS}\sum_{j} iS^{z}_{j} \\
&= 1
\end{align*}

（译注：照录原书。首行第一个括号中的 $S^{j}_{\lambda}$ 显系 $S^{y}_{\lambda}$ 之误。又，按对易关系 $[S^{x},S^{y}] = iS^{z}$，第二行左端 $-i(S^{x}_{\lambda}S^{y}_{\lambda} - S^{y}_{\lambda}S^{x}_{\lambda})$ 应等于 $\frac{1}{NS}\sum_{j} S^{z}_{j}$（基态中 $\sum_{j}S^{z}_{j} = NS$），故该行系数中的两个 $i$ 多出一个，否则与末行 “$= 1$” 不相容。）

Klein 与 Smith (99) 最先指出，情况之所以得救，是由于与量子自旋矢量 $\mathbf{S}$ 相对应的经典自旋是 $S^{'} = \sqrt{S(S+1)}$，也就是说，$S(S+1) = S_{xj}^{2} + S_{yj}^{2} + S_{zj}^{2}$。于是，$S_{zj}^{2} = S^{2} - (S_{xj}^{2} + S_{yj}^{2} - S)$；再次应用近似对易关系

\begin{equation*}
S_{xj}S_{yj} - S_{yj}S_{xj} = iS_{zj} \simeq iS
\end{equation*}

就得到

\begin{align*}
S_{zj}^{2} &= S^{2} - \left[\frac{1}{2}\left(S^{+}_{j}S^{-}_{j} + S^{-}_{j}S^{+}_{j}\right) - S\right] \\
&\simeq S^{2} - S^{-}S^{+}
\end{align*}

因此，只有把 $S^{2}$ 改成 $S(S+1)$，(26) 才是正确的对应原理哈密顿量；而这样一改，它便与量子结果完全一致：

\begin{equation*}
\mathcal{H} = -NZJS^{2} + \sum_{\lambda} \omega_{\lambda}\, S^{-}_{\lambda}S^{+}_{\lambda}
\end{equation*}

其结果是经典理论与量子理论之间的完全对应。同激子情形中我们略去“逆向图”（backward diagrams）时一样，零点运动效应不过是一种相当平凡的数学练习，只与单个原子的内部动力学有关；耦合项的零点平均值净效果为零，因为零点效应要对所有自旋波 $\lambda$ 求和，而 $\sum_{\lambda}\cos\lambda = 0$（在这里，$\cos\lambda$ 就是耦合项的 Fourier 变换）。

还有第二种用途非常普遍的处理磁性问题的半经典方法，它是由 Landau、Lifshitz、Kittel 与 Herring (100) 提出的。这一方法非常接近于声子的 Debye 理论。其想法是立足于相当粗粒的观点，把磁化强度看作一个连续的磁化场 $\mathbf{M}(r)$，它服从连续量子场的对易关系

\begin{equation*}
\mathbf{M}(r)\times\mathbf{M}(r') = i\mathbf{M}(r)\,\delta(r-r')\ \left(\times\ \text{consts.}\right)
\end{equation*}

哈密顿量依赖于 $(S_{j} - S_{k})^{2}$ 这一形式提示我们，作为半经验的哈密顿量密度可以取

\begin{equation*}
\mathcal{H} = A\int dr\left(\nabla\mathbf{M}(r)\right)^{2} \tag{27}
\end{equation*}

可以证明，对绝缘铁磁体这一特殊情形它是正确的；而且显而易见，它也是任何短程、各向同性的力的宏观必然后果——正如弹性理论必然给出一个形变 $(\nabla u)$ 的二次型哈密顿量一样（$u$ 为位移）。由此导出了普遍的磁化运动方程，人们可以猜想，它们在所有铁磁体中都成立：

\begin{equation*}
\frac{d\mathbf{M}}{dt} = A\left(\mathbf{M}\times\nabla^{2}\mathbf{M}\right) + \gamma\left(\mathbf{M}\times\mathbf{H}_{\mathrm{ext}}\right) + \left(\text{耗散项}\right)
\end{equation*}

这就是 Landau-Lifshitz 方程，它在所有自旋波理论和铁磁共振理论中都起着基础性的作用。

像 (27) 这样的哈密顿量——以及弹性理论中相应的哈密顿量

\begin{equation*}
\mathcal{H} = \int \nabla\mathbf{u}:\mathbf{C}:\nabla\mathbf{u} \qquad \left(\mathbf{C}\ \text{为刚度弹性张量}\right)
\end{equation*}

——所共有的一个有趣性质，就是下述普遍论证：具有短程力的二维系统不可能有稳定的长程序。我们很容易相信，系统的长波长运动——至少当其频率很低时——在任一有限温度下都按经典方式行为。于是按能量均分，每个自由度 $k$ 分得 $k_{\mathrm{B}}T$ 的能量，而能量是

\begin{equation*}
\propto k^{2}\left(\nabla\mathbf{M}(k)\right)^{2} \qquad \text{故} \qquad \left\langle\left(\nabla\mathbf{M}(k)\right)^{2}\right\rangle = \frac{k_{\mathrm{B}}T}{k^{2}}
\end{equation*}

于是涨落积分

\begin{equation*}
\left\langle\Delta M^{2}\right\rangle = \int k^{D-1}\,dk\,\frac{k_{\mathrm{B}}T}{k^{2}} \qquad \left(D\ \text{为维数}\right)
\end{equation*}

在磁性与弹性两种情形中都在 $D = 2$ 时发散。换一种说法就是：在二维情形下，人们可以造就一个尺寸为

% ——上句在原书 p.174 末中断于 “of size”，PDF p190 起由后续代理续译衔接。
